\documentclass{amsart}
% For dealing with references we use the comment environment
\usepackage{verbatim}
\newenvironment{reference}{\comment}{\endcomment}
%\newenvironment{reference}{}{}
\newenvironment{slogan}{\comment}{\endcomment}
\newenvironment{history}{\comment}{\endcomment}

% For commutative diagrams we use Xy-pic
\usepackage[all]{xy}

% We use 2cell for 2-commutative diagrams.
\xyoption{2cell}
\UseAllTwocells

% We use multicol for the list of chapters between chapters
\usepackage{multicol}

% This is generally recommended for better output
\usepackage{lmodern}
\usepackage[T1]{fontenc}

% For cross-file-references

% Package for hypertext links:
\usepackage{hyperref}

% For any local file, say "hello.tex" you want to link to please

% Theorem environments.
%
\theoremstyle{plain}
\newtheorem{theorem}[subsection]{Theorem}
\newtheorem{proposition}[subsection]{Proposition}
\newtheorem{lemma}[subsection]{Lemma}

\theoremstyle{definition}
\newtheorem{definition}[subsection]{Definition}
\newtheorem{example}[subsection]{Example}
\newtheorem{exercise}[subsection]{Exercise}
\newtheorem{situation}[subsection]{Situation}

\theoremstyle{remark}
\newtheorem{remark}[subsection]{Remark}
\newtheorem{remarks}[subsection]{Remarks}

\numberwithin{equation}{subsection}

% Macros
%
\def\lim{\mathop{\mathrm{lim}}\nolimits}
\def\colim{\mathop{\mathrm{colim}}\nolimits}
\def\Spec{\mathop{\mathrm{Spec}}}
\def\Hom{\mathop{\mathrm{Hom}}\nolimits}
\def\Ext{\mathop{\mathrm{Ext}}\nolimits}
\def\SheafHom{\mathop{\mathcal{H}\!\mathit{om}}\nolimits}
\def\SheafExt{\mathop{\mathcal{E}\!\mathit{xt}}\nolimits}
\def\Sch{\mathit{Sch}}
\def\Mor{\mathop{\mathrm{Mor}}\nolimits}
\def\Ob{\mathop{\mathrm{Ob}}\nolimits}
\def\Sh{\mathop{\mathit{Sh}}\nolimits}
\def\NL{\mathop{N\!L}\nolimits}
\def\CH{\mathop{\mathrm{CH}}\nolimits}
\def\proetale{{pro\text{-}\acute{e}tale}}
\def\etale{{\acute{e}tale}}
\def\QCoh{\mathit{QCoh}}
\def\Ker{\mathop{\mathrm{Ker}}}
\def\Im{\mathop{\mathrm{Im}}}
\def\Coker{\mathop{\mathrm{Coker}}}
\def\Coim{\mathop{\mathrm{Coim}}}

% Boxtimes
%
\DeclareMathSymbol{\boxtimes}{\mathbin}{AMSa}{"02}

%
% Macros for moduli stacks/spaces
%
\def\QCohstack{\mathcal{QC}\!\mathit{oh}}
\def\Cohstack{\mathcal{C}\!\mathit{oh}}
\def\Spacesstack{\mathcal{S}\!\mathit{paces}}
\def\Quotfunctor{\mathrm{Quot}}
\def\Hilbfunctor{\mathrm{Hilb}}
\def\Curvesstack{\mathcal{C}\!\mathit{urves}}
\def\Polarizedstack{\mathcal{P}\!\mathit{olarized}}
\def\Complexesstack{\mathcal{C}\!\mathit{omplexes}}
% \Pic is the operator that assigns to X its picard group, usage \Pic(X)
% \Picardstack_{X/B} denotes the Picard stack of X over B
% \Picardfunctor_{X/B} denotes the Picard functor of X over B
\def\Pic{\mathop{\mathrm{Pic}}\nolimits}
\def\Picardstack{\mathcal{P}\!\mathit{ic}}
\def\Picardfunctor{\mathrm{Pic}}
\def\Deformationcategory{\mathcal{D}\!\mathit{ef}}

\setlength{\emergencystretch}{3em}
\begin{document}
\title[Stacks Project, Tag 0885]{376 --- Algebraization of coherent formal systems on projective schemes over complete Noetherian rings}
\maketitle
\noindent Copyright (C) 2005--2025 Johan de Jong. Stacks Project authors. Source: \href{https://stacks.math.columbia.edu/tag/0885}{Tag 0885}.

\noindent Editorial difficulty note (not author text). Difficulty H2: Essential surjectivity requires constructing a coherent algebraic presentation for a compatible formal system. Uniform twisting vanishing lifts finitely many generators through every infinitesimal level, producing a surjection from a completed finite sum of ample twists; a second presentation and full faithfulness then algebraize the cokernel..

\noindent Editorial provenance note (not author text). History: excerpt from revision \texttt{a0444\allowbreak 6e57e\allowbreak c1fbc\allowbreak 252a8\allowbreak 71afc\allowbreak ec775\allowbreak 2fb28\allowbreak 07b14}, coherent.tex, lines 6811--6867. Reference presentation was adapted; original-author context and core support blocks were retained or added. The mathematical wording is unchanged. Licensed under GNU FDL 1.2 or later; no invariant sections or cover texts. See ../licenses/COPYING. Context blocks reproduce author definitions and supporting results; they are not additional counted problems.

\section{Coherent formal modules}
\label{section-coherent-formal}

\noindent
As we do not yet have the theory of formal schemes to our disposal, we
develop a bit of language that replaces the notion of a ``coherent module on a
Noetherian adic formal scheme''.

\medskip\noindent
Let $X$ be a Noetherian scheme and let $\mathcal{I} \subset \mathcal{O}_X$
be a quasi-coherent sheaf of ideals. We will consider inverse
systems $(\mathcal{F}_n)$ of coherent $\mathcal{O}_X$-modules such that
\begin{enumerate}
\item $\mathcal{F}_n$ is annihilated by $\mathcal{I}^n$, and
\item the transition maps induce isomorphisms
$\mathcal{F}_{n + 1}/\mathcal{I}^n\mathcal{F}_{n + 1} \to \mathcal{F}_n$.
\end{enumerate}
A morphism of such inverse systems is defined as usual.
Let us denote the category of these inverse systems with
$\textit{Coh}(X, \mathcal{I})$. We are going to proceed by proving
a bunch of lemmas about objects in this category. In fact, most
of the lemmas that follow are straightforward consequences of the following
description of the category in the affine case.



\noindent
Let $X$ be a Noetherian scheme and let $\mathcal{I} \subset \mathcal{O}_X$
be a quasi-coherent sheaf of ideals. There is a functor
\begin{equation}
\label{equation-completion-functor}
\textit{Coh}(\mathcal{O}_X) \longrightarrow \textit{Coh}(X, \mathcal{I}), \quad
\mathcal{F} \longmapsto \mathcal{F}^\wedge
\end{equation}
which associates to the coherent $\mathcal{O}_X$-module $\mathcal{F}$
the object $\mathcal{F}^\wedge = (\mathcal{F}/\mathcal{I}^n\mathcal{F})$
of $\textit{Coh}(X, \mathcal{I})$.

\begin{theorem}[Theorem on formal functions]
\label{theorem-formal-functions}
Let $A$ be a Noetherian ring.
Let $I \subset A$ be an ideal.
Let $f : X \to \Spec(A)$ be a proper morphism.
Let $\mathcal{F}$ be a coherent sheaf on $X$.
Fix $p \geq 0$.
The system of maps
$$
H^p(X, \mathcal{F})/I^nH^p(X, \mathcal{F})
\longrightarrow
H^p(X, \mathcal{F}/I^n\mathcal{F})
$$
define an isomorphism of limits
$$
H^p(X, \mathcal{F})^\wedge
\longrightarrow
\lim_n H^p(X, \mathcal{F}/I^n\mathcal{F})
$$
where the left hand side is the completion of the $A$-module
$H^p(X, \mathcal{F})$ with respect to the ideal $I$, see
Algebra, Section \href{https://stacks.math.columbia.edu/tag/00M9}{section completion (Tag 00M9)}.
Moreover, this is in fact a homeomorphism for the limit topologies.
\end{theorem}

\begin{lemma}
\label{lemma-vanishing-projective}
Let $A$ be Noetherian ring and $I \subset A$ an ideal.
Let $f : X \to \Spec(A)$ be a proper morphism and let
$\mathcal{L}$ be an $f$-ample invertible sheaf. Let
$\mathcal{I} = I\mathcal{O}_X$. Let $(\mathcal{F}_n)$ be an
object of $\textit{Coh}(X, \mathcal{I})$. Then there exists an
integer $d_0$ such that
$$
H^1(X, \Ker(\mathcal{F}_{n + 1} \to \mathcal{F}_n)
\otimes \mathcal{L}^{\otimes d} )
= 0
$$
for all $n \geq 0$ and all $d \geq d_0$.
\end{lemma}

\begin{proof}
Set $B = \bigoplus I^n/I^{n + 1}$ and
$\mathcal{B} = \bigoplus \mathcal{I}^n/\mathcal{I}^{n + 1} = f^*\widetilde{B}$.
By Lemma \ref{lemma-finite-over-rees-algebra} the graded quasi-coherent
$\mathcal{B}$-module
$\mathcal{G} = \bigoplus \Ker(\mathcal{F}_{n + 1} \to \mathcal{F}_n)$
is of finite type. Hence the lemma follows from
Lemma \ref{lemma-graded-finiteness} part (2).
\end{proof}


\begin{lemma}
\label{lemma-finite-over-rees-algebra}
Let $X$ be a Noetherian scheme and let $\mathcal{I} \subset \mathcal{O}_X$
be a quasi-coherent sheaf of ideals. If $(\mathcal{F}_n)$ is an object of
$\textit{Coh}(X, \mathcal{I})$ then
$\bigoplus \Ker(\mathcal{F}_{n + 1} \to \mathcal{F}_n)$ is
a finite type, graded, quasi-coherent
$\bigoplus \mathcal{I}^n/\mathcal{I}^{n + 1}$-module.
\end{lemma}

\begin{proof}
The question is local on $X$ hence we may assume $X$ is affine, i.e.,
we have a situation as in Lemma \ref{lemma-inverse-systems-affine}.
In this case, if $(\mathcal{F}_n)$ corresponds to the finite $A^\wedge$
module $M$, then $\bigoplus \Ker(\mathcal{F}_{n + 1} \to \mathcal{F}_n)$
corresponds to $\bigoplus I^nM/I^{n + 1}M$ which is clearly a finite
module over $\bigoplus I^n/I^{n + 1}$.
\end{proof}


\begin{lemma}
\label{lemma-graded-finiteness}
Let $A$ be a Noetherian ring.
Let $B$ be a finitely generated graded $A$-algebra.
Let $f : X \to \Spec(A)$ be a proper morphism.
Set $\mathcal{B} = f^*\widetilde B$.
Let $\mathcal{F}$ be a quasi-coherent
graded $\mathcal{B}$-module of finite type.
\begin{enumerate}
\item For every $p \geq 0$ the graded $B$-module $H^p(X, \mathcal{F})$
is a finite $B$-module.
\item If $\mathcal{L}$ is an ample invertible $\mathcal{O}_X$-module,
then there exists an integer $d_0$ such that
$H^p(X, \mathcal{F} \otimes \mathcal{L}^{\otimes d}) = 0$
for all $p > 0$ and $d \geq d_0$.
\end{enumerate}
\end{lemma}

\begin{proof}
To prove this we consider the fibre product diagram
$$
\xymatrix{
X' = \Spec(B) \times_{\Spec(A)} X
\ar[r]_-\pi \ar[d]_{f'} &
X \ar[d]^f \\
\Spec(B) \ar[r] &
\Spec(A)
}
$$
Note that $f'$ is a proper morphism, see
Morphisms, Lemma \href{https://stacks.math.columbia.edu/tag/01W4}{lemma base change proper (Tag 01W4)}.
Also, $B$ is a finitely generated $A$-algebra, and hence
Noetherian (Algebra, Lemma \href{https://stacks.math.columbia.edu/tag/00FN}{lemma Noetherian permanence (Tag 00FN)}).
This implies that $X'$ is a Noetherian scheme
(Morphisms, Lemma \href{https://stacks.math.columbia.edu/tag/01T6}{lemma finite type noetherian (Tag 01T6)}).
Note that $X'$ is the relative spectrum of the quasi-coherent
$\mathcal{O}_X$-algebra $\mathcal{B}$ by
Constructions, Lemma \href{https://stacks.math.columbia.edu/tag/01LX}{lemma spec properties (Tag 01LX)}.
Since $\mathcal{F}$ is a quasi-coherent $\mathcal{B}$-module
we see that there is a unique quasi-coherent
$\mathcal{O}_{X'}$-module $\mathcal{F}'$ such that
$\pi_*\mathcal{F}' = \mathcal{F}$, see
Morphisms, Lemma \href{https://stacks.math.columbia.edu/tag/01SB}{lemma affine equivalence modules (Tag 01SB)}
Since $\mathcal{F}$ is finite type as a $\mathcal{B}$-module we
conclude that $\mathcal{F}'$ is a finite type
$\mathcal{O}_{X'}$-module (details omitted). In other words,
$\mathcal{F}'$ is a coherent $\mathcal{O}_{X'}$-module
(Lemma \href{https://stacks.math.columbia.edu/tag/01XZ}{lemma coherent Noetherian (Tag 01XZ)}).
Since the morphism $\pi : X' \to X$ is affine we have
$$
H^p(X, \mathcal{F}) = H^p(X', \mathcal{F}')
$$
by Lemma \href{https://stacks.math.columbia.edu/tag/089W}{lemma relative affine cohomology (Tag 089W)}.
Thus (1) follows from
Lemma \href{https://stacks.math.columbia.edu/tag/02O6}{lemma proper over affine cohomology finite (Tag 02O6)}.
Given $\mathcal{L}$ as in (2) we set
$\mathcal{L}' = \pi^*\mathcal{L}$. Note that $\mathcal{L}'$ is
ample on $X'$ by
Morphisms, Lemma \href{https://stacks.math.columbia.edu/tag/0892}{lemma pullback ample tensor relatively ample (Tag 0892)}.
By the projection formula
(Cohomology, Lemma \href{https://stacks.math.columbia.edu/tag/01E8}{lemma projection formula (Tag 01E8)}) we have
$\pi_*(\mathcal{F}' \otimes \mathcal{L}') = \mathcal{F} \otimes \mathcal{L}$.
Thus part (2) follows by the same reasoning as above from
Lemma \href{https://stacks.math.columbia.edu/tag/02O1}{lemma kill by twisting (Tag 02O1)}.
\end{proof}


\begin{lemma}
\label{lemma-inverse-systems-affine}
If $X = \Spec(A)$ is the spectrum of a Noetherian ring and
$\mathcal{I}$ is the quasi-coherent sheaf of ideals associated to the ideal
$I \subset A$, then $\textit{Coh}(X, \mathcal{I})$ is equivalent to the
category of finite $A^\wedge$-modules where $A^\wedge$ is the completion
of $A$ with respect to $I$.
\end{lemma}

\begin{proof}
Let $\text{Mod}^{fg}_{A, I}$ be the category of inverse systems $(M_n)$
of finite $A$-modules satisfying: (1) $M_n$ is annihilated by $I^n$ and (2)
$M_{n + 1}/I^nM_{n + 1} = M_n$. By the correspondence between coherent
sheaves on $X$ and finite $A$-modules (Lemma \href{https://stacks.math.columbia.edu/tag/01XZ}{lemma coherent Noetherian (Tag 01XZ)})
it suffices to show $\text{Mod}^{fg}_{A, I}$ is equivalent to the category of
finite $A^\wedge$-modules. To see this it suffices to prove that given
an object $(M_n)$ of $\text{Mod}^{fg}_{A, I}$ the module
$$
M = \lim M_n
$$
is a finite $A^\wedge$-module and that $M/I^nM = M_n$. As the transition
maps are surjective, we see that $M \to M_1$ is surjective.
Pick $x_1, \ldots, x_t \in M$ which map to generators of $M_1$.
This induces a map of systems $(A/I^n)^{\oplus t} \to M_n$.
By Nakayama's lemma (Algebra, Lemma \ref{algebra-lemma-NAK}) these maps are
surjective. Let $K_n \subset (A/I^n)^{\oplus t}$ be the kernel.
Property (2) implies that $K_{n + 1} \to K_n$ is surjective, in particular
the system $(K_n)$ satisfies the Mittag-Leffler condition.
By Homology, Lemma \href{https://stacks.math.columbia.edu/tag/02N1}{lemma Mittag Leffler (Tag 02N1)}
we obtain an exact sequence
$0 \to K \to (A^\wedge)^{\oplus t} \to M \to 0$
with $K = \lim K_n$.
Hence $M$ is a finite $A^\wedge$-module.
As $K \to K_n$ is surjective it follows that
$$
M/I^nM = \Coker(K \to (A/I^n)^{\oplus t}) =
(A/I^n)^{\oplus t}/K_n = M_n
$$
as desired.
\end{proof}


\begin{lemma}
\label{lemma-inverse-systems-abelian}
Let $X$ be a Noetherian scheme and let $\mathcal{I} \subset \mathcal{O}_X$
be a quasi-coherent sheaf of ideals.
\begin{enumerate}
\item The category $\textit{Coh}(X, \mathcal{I})$ is abelian.
\item For $U \subset X$ open the restriction functor
$\textit{Coh}(X, \mathcal{I}) \to \textit{Coh}(U, \mathcal{I}|_U)$
is exact.
\item Exactness in $\textit{Coh}(X, \mathcal{I})$ may be checked by
restricting to the members of an open covering of $X$.
\end{enumerate}
\end{lemma}

\begin{proof}
Let $\alpha =(\alpha_n) : (\mathcal{F}_n) \to (\mathcal{G}_n)$ be a morphism of
$\textit{Coh}(X, \mathcal{I})$. The cokernel of $\alpha$ is the inverse system
$(\Coker(\alpha_n))$ (details omitted). To describe the kernel let
$$
\mathcal{K}'_{l, m} = \Im(\Ker(\alpha_l) \to \mathcal{F}_m)
$$
for $l \geq m$.
We claim:
\begin{enumerate}
\item[(a)] the inverse system $(\mathcal{K}'_{l, m})_{l \geq m}$ is
eventually constant, say with value $\mathcal{K}'_m$,
\item[(b)] the system $(\mathcal{K}'_m/\mathcal{I}^n\mathcal{K}'_m)_{m \geq n}$
is eventually constant, say with value $\mathcal{K}_n$,
\item[(c)] the system $(\mathcal{K}_n)$ forms an object of
$\textit{Coh}(X, \mathcal{I})$, and
\item[(d)] this object is the kernel of $\alpha$.
\end{enumerate}
To see (a), (b), and (c) we may work affine locally, say $X = \Spec(A)$
and $\mathcal{I}$ corresponds to the ideal $I \subset A$. By
Lemma \ref{lemma-inverse-systems-affine}
$\alpha$ corresponds to a map $f : M \to N$ of finite $A^\wedge$-modules.
Denote $K = \Ker(f)$. Note that $A^\wedge$ is a Noetherian
ring (Algebra, Lemma \href{https://stacks.math.columbia.edu/tag/0316}{lemma completion Noetherian Noetherian (Tag 0316)}).
Choose an integer $c \geq 0$ such that
$K \cap I^n M \subset I^{n - c}K$ for $n \geq c$
(Algebra, Lemma \href{https://stacks.math.columbia.edu/tag/00IN}{lemma Artin Rees (Tag 00IN)})
and which satisfies Algebra, Lemma \href{https://stacks.math.columbia.edu/tag/00IO}{lemma map AR (Tag 00IO)}
for the map $f$ and the ideal $I^\wedge = IA^\wedge$. Then
$\mathcal{K}'_{l, m}$ corresponds to the $A$-module
$$
K'_{l, m} = \frac{a^{-1}(I^lN) + I^mM}{I^mM} =
\frac{K + I^{l - c}f^{-1}(I^cN) + I^mM}{I^mM} =
\frac{K + I^mM}{I^mM}
$$
where the last equality holds if $l \geq m + c$. So $\mathcal{K}'_m$
corresponds to the $A$-module $K/K \cap I^mM$ and
$\mathcal{K}'_m/\mathcal{I}^n\mathcal{K}'_m$ corresponds to
$$
\frac{K}{K \cap I^mM + I^nK} = \frac{K}{I^nK}
$$
for $m \geq n + c$ by our choice of $c$ above. Hence $\mathcal{K}_n$
corresponds to $K/I^nK$.

\medskip\noindent
We prove (d). It is clear from the description on affines above that
the composition $(\mathcal{K}_n) \to (\mathcal{F}_n) \to (\mathcal{G}_n)$
is zero. Let $\beta : (\mathcal{H}_n) \to (\mathcal{F}_n)$
be a morphism such that $\alpha \circ \beta = 0$. Then
$\mathcal{H}_l \to \mathcal{F}_l$ maps into $\Ker(\alpha_l)$.
Since $\mathcal{H}_m = \mathcal{H}_l/\mathcal{I}^m\mathcal{H}_l$
for $l \geq m$ we obtain a system of maps
$\mathcal{H}_m \to \mathcal{K}'_{l, m}$. Thus a map
$\mathcal{H}_m \to \mathcal{K}_m'$. Since
$\mathcal{H}_n = \mathcal{H}_m/\mathcal{I}^n\mathcal{H}_m$ we obtain
a system of maps $\mathcal{H}_n \to \mathcal{K}'_m/\mathcal{I}^n\mathcal{K}'_m$
and hence a map $\mathcal{H}_n \to \mathcal{K}_n$ as desired.

\medskip\noindent
To finish the proof of (1) we still have to show that $\Coim = \Im$
in $\textit{Coh}(X, \mathcal{I})$. We have seen above that taking
kernels and cokernels commutes, over affines, with the description
of $\textit{Coh}(X, \mathcal{I})$ as a category of modules. Since
$\Im = \Coim$ holds in the category of modules
this gives $\Coim = \Im$ in $\textit{Coh}(X, \mathcal{I})$.
Parts (2) and (3) of the lemma are immediate from our construction
of kernels and cokernels.
\end{proof}


\begin{lemma}
\label{lemma-inverse-systems-surjective}
Let $X$ be a Noetherian scheme and let $\mathcal{I} \subset \mathcal{O}_X$
be a quasi-coherent sheaf of ideals. A map
$(\mathcal{F}_n) \to (\mathcal{G}_n)$ is surjective in
$\textit{Coh}(X, \mathcal{I})$
if and only if $\mathcal{F}_1 \to \mathcal{G}_1$ is surjective.
\end{lemma}

\begin{proof}
Omitted. Hint: Look on affine opens, use
Lemma \ref{lemma-inverse-systems-affine}, and use
Algebra, Lemma \ref{algebra-lemma-NAK}.
\end{proof}


\begin{lemma}
\label{lemma-exact}
The functor (\ref{equation-completion-functor}) is exact.
\end{lemma}

\begin{proof}
It suffices to check this locally on $X$. Hence we may assume $X$ is
affine, i.e., we have a situation as in
Lemma \ref{lemma-inverse-systems-affine}.
The functor is the functor $\text{Mod}^{fg}_A \to \text{Mod}^{fg}_{A^\wedge}$
which associates to a finite $A$-module $M$ the completion $M^\wedge$.
Thus the result follows from
Algebra, Lemma \href{https://stacks.math.columbia.edu/tag/00MB}{lemma completion flat (Tag 00MB)}.
\end{proof}


\begin{lemma}
\label{lemma-completion-internal-hom}
Let $X$ be a Noetherian scheme and let $\mathcal{I} \subset \mathcal{O}_X$
be a quasi-coherent sheaf of ideals. Let $\mathcal{F}$, $\mathcal{G}$ be
coherent $\mathcal{O}_X$-modules. Set
$\mathcal{H} = \SheafHom_{\mathcal{O}_X}(\mathcal{G}, \mathcal{F})$.
Then
$$
\lim H^0(X, \mathcal{H}/\mathcal{I}^n\mathcal{H}) =
\Mor_{\textit{Coh}(X, \mathcal{I})}
(\mathcal{G}^\wedge, \mathcal{F}^\wedge).
$$
\end{lemma}

\begin{proof}
To prove this we may work affine locally on $X$.
Hence we may assume $X = \Spec(A)$ and $\mathcal{F}$, $\mathcal{G}$
given by finite $A$-module $M$ and $N$. Then $\mathcal{H}$
corresponds to the finite $A$-module $H = \Hom_A(M, N)$.
The statement of the lemma becomes the statement
$$
H^\wedge = \Hom_{A^\wedge}(M^\wedge, N^\wedge)
$$
via the equivalence of Lemma \ref{lemma-inverse-systems-affine}.
By Algebra, Lemma \href{https://stacks.math.columbia.edu/tag/00MB}{lemma completion flat (Tag 00MB)}
(used 3 times) we have
$$
H^\wedge = \Hom_A(M, N) \otimes_A A^\wedge =
\Hom_{A^\wedge}(M \otimes_A A^\wedge, N \otimes_A A^\wedge) =
\Hom_{A^\wedge}(M^\wedge, N^\wedge)
$$
where the second equality uses that $A^\wedge$ is flat over $A$
(see More on Algebra, Lemma
\href{https://stacks.math.columbia.edu/tag/087R}{lemma pseudo coherence and base change ext (Tag 087R)}).
The lemma follows.
\end{proof}


\begin{lemma}
\label{lemma-fully-faithful}
Let $A$ be Noetherian ring complete with respect to an ideal $I$.
Let $f : X \to \Spec(A)$ be a proper morphism. Let
$\mathcal{I} = I\mathcal{O}_X$.
Then the functor (\ref{equation-completion-functor}) is fully faithful.
\end{lemma}

\begin{proof}
Let $\mathcal{F}$, $\mathcal{G}$ be coherent $\mathcal{O}_X$-modules.
Then $\mathcal{H} = \SheafHom_{\mathcal{O}_X}(\mathcal{G}, \mathcal{F})$
is a coherent $\mathcal{O}_X$-module, see
Modules, Lemma \href{https://stacks.math.columbia.edu/tag/01CQ}{lemma internal hom locally kernel direct sum (Tag 01CQ)}.
By Lemma \ref{lemma-completion-internal-hom} the map
$$
\lim_n H^0(X, \mathcal{H}/\mathcal{I}^n\mathcal{H})
\to
\Mor_{\textit{Coh}(X, \mathcal{I})}
(\mathcal{G}^\wedge, \mathcal{F}^\wedge)
$$
is bijective. Hence fully faithfulness of
(\ref{equation-completion-functor}) follows from the theorem on formal
functions (Lemma \ref{lemma-spell-out-theorem-formal-functions})
for the coherent sheaf $\mathcal{H}$.
\end{proof}


\begin{lemma}
\label{lemma-spell-out-theorem-formal-functions}
Let $A$ be a ring. Let $I \subset A$ be an ideal. Assume $A$ is
Noetherian and complete with respect to $I$.
Let $f : X \to \Spec(A)$ be a proper morphism.
Let $\mathcal{F}$ be a coherent sheaf on $X$.
Then
$$
H^p(X, \mathcal{F}) = \lim_n H^p(X, \mathcal{F}/I^n\mathcal{F})
$$
for all $p \geq 0$.
\end{lemma}

\begin{proof}
This is a reformulation of the theorem on formal functions
(Theorem \ref{theorem-formal-functions}) in the
case of a complete Noetherian base ring. Namely, in this case the
$A$-module $H^p(X, \mathcal{F})$ is finite
(Lemma \href{https://stacks.math.columbia.edu/tag/02O6}{lemma proper over affine cohomology finite (Tag 02O6)}) hence
$I$-adically complete (Algebra, Lemma \href{https://stacks.math.columbia.edu/tag/00MA}{lemma completion tensor (Tag 00MA)})
and we see that completion on the left hand side is not necessary.
\end{proof}


\begin{lemma}[Nakayama's lemma]
\label{algebra-lemma-NAK}
\begin{reference}
\href{https://stacks.math.columbia.edu/bibliography}{Bibliography: MatCA (1.M Lemma (NAK) page 11)}
\end{reference}
\begin{history}
We quote from \href{https://stacks.math.columbia.edu/bibliography}{Bibliography: MatCA}: ``This simple but
important lemma is due to T.~Nakayama, G.~Azumaya and W.~Krull. Priority
is obscure, and although it is usually called the Lemma of Nakayama, late
Prof.~Nakayama did not like the name.''
\end{history}
Let $R$ be a ring with Jacobson radical $\text{rad}(R)$.
Let $M$ be an $R$-module. Let $I \subset R$
be an ideal.
\begin{enumerate}
\item
\label{algebra-item-nakayama}
If $IM = M$ and $M$ is finite, then there exists an $f \in 1 + I$ such that
$fM = 0$.
\item If $IM = M$, $M$ is finite, and $I \subset \text{rad}(R)$, then $M = 0$.
\item If $N, N' \subset M$, $M = N + IN'$, and $N'$ is finite,
then there exists an $f \in 1 + I$ such that $fM \subset N$ and $M_f = N_f$.
\item If $N, N' \subset M$, $M = N + IN'$, $N'$ is finite, and
$I \subset \text{rad}(R)$, then $M = N$.
\item If $N \to M$ is a module map, $N/IN \to M/IM$ is
surjective, and $M$ is finite, then there exists an $f \in 1 + I$
such that $N_f \to M_f$ is surjective.
\item If $N \to M$ is a module map, $N/IN \to M/IM$ is
surjective, $M$ is finite, and $I \subset \text{rad}(R)$,
then $N \to M$ is surjective.
\item If $x_1, \ldots, x_n \in M$ generate $M/IM$ and $M$ is finite,
then there exists an $f \in 1 + I$ such that $x_1, \ldots, x_n$
generate $M_f$ over $R_f$.
\item If $x_1, \ldots, x_n \in M$ generate $M/IM$, $M$ is finite, and
$I \subset \text{rad}(R)$, then $M$ is generated by $x_1, \ldots, x_n$.
\item If $IM = M$, $I$ is nilpotent, then $M = 0$.
\item If $N, N' \subset M$, $M = N + IN'$, and $I$ is nilpotent then $M = N$.
\item If $N \to M$ is a module map, $I$ is nilpotent, and $N/IN \to M/IM$
is surjective, then $N \to M$ is surjective.
\item If $\{x_\alpha\}_{\alpha \in A}$ is a set of elements of $M$
which generate $M/IM$ and $I$ is nilpotent, then $M$ is generated
by the $x_\alpha$.
\end{enumerate}
\end{lemma}

\begin{proof}
Proof of (\ref{algebra-item-nakayama}). Choose generators $y_1, \ldots, y_m$ of $M$
over $R$. For each $i$ we can write $y_i = \sum z_{ij} y_j$ with
$z_{ij} \in I$ (since $M = IM$).
In other words $\sum_j (\delta_{ij} - z_{ij})y_j = 0$.
Let $f$ be the determinant of the $m \times m$ matrix
$A = (\delta_{ij} - z_{ij})$. Note that $f \in 1 + I$
(since the matrix $A$ is entrywise congruent to the
$m \times m$ identity matrix modulo $I$).
By Lemma \href{https://stacks.math.columbia.edu/tag/07DQ}{lemma matrix left inverse (Tag 07DQ)} (1),
there exists an $m \times m$
matrix $B$ such that $BA = f 1_{m \times m}$. Writing out we see that
$\sum_{i} b_{hi} a_{ij} = f \delta_{hj}$ for all
$h$ and $j$; hence, $\sum_{i, j} b_{hi} a_{ij} y_j
= \sum_{j} f \delta_{hj} y_j = f y_h$ for every $h$.
In other words, $0 = f y_h$ for every $h$ (since each
$i$ satisfies $\sum_j a_{ij} y_j = 0$).
This implies that $f$ annihilates $M$.

\medskip\noindent
By Lemma \href{https://stacks.math.columbia.edu/tag/0AME}{lemma contained in radical (Tag 0AME)} an element of $1 + \text{rad}(R)$ is
invertible element of $R$. Hence we see that (\ref{algebra-item-nakayama}) implies
(2). We obtain (3) by applying (1) to $M/N$ which is finite as $N'$ is finite.
We obtain (4) by applying (2) to $M/N$ which is finite as $N'$ is finite.
We obtain (5) by applying (3) to $M$ and the submodules $\Im(N \to M)$
and $M$. We obtain (6) by applying (4) to $M$ and the submodules
$\Im(N \to M)$ and $M$.
We obtain (7) by applying (5) to the map $R^{\oplus n} \to M$,
$(a_1, \ldots, a_n) \mapsto a_1x_1 + \ldots + a_nx_n$.
We obtain (8) by applying (6) to the map $R^{\oplus n} \to M$,
$(a_1, \ldots, a_n) \mapsto a_1x_1 + \ldots + a_nx_n$.

\medskip\noindent
Part (9) holds because if $M = IM$ then $M = I^nM$ for all $n \geq 0$
and $I$ being nilpotent means $I^n = 0$ for some $n \gg 0$. Parts
(10), (11), and (12) follow from (9) by the arguments used above.
\end{proof}


\begin{lemma}
\label{algebra-lemma-radical-completion}
Let $R$ be a ring, let $I \subset R$ be an ideal, and let
$R^\wedge = \lim R/I^n$.
\begin{enumerate}
\item any element of $R^\wedge$ which maps to a unit of $R/I$ is a unit,
\item any element of $1 + I$ maps to an invertible element of $R^\wedge$,
\item any element of $1 + IR^\wedge$ is invertible in $R^\wedge$, and
\item the ideals $IR^\wedge$ and $\Ker(R^\wedge \to R/I)$ are contained
in the Jacobson radical of $R^\wedge$.
\end{enumerate}
\end{lemma}

\begin{proof}
Let $x \in R^\wedge$ map to a unit $x_1$ in $R/I$.
Then $x$ maps to a unit $x_n$ in $R/I^n$ for every $n$ by
Lemma \href{https://stacks.math.columbia.edu/tag/0AMG}{lemma locally nilpotent unit (Tag 0AMG)}.
Hence $y = (x_n^{-1}) \in \lim R/I^n = R^\wedge$ is an inverse to $x$.
Parts (2) and (3) follow immediately from (1).
Part (4) follows from (1) and Lemma \href{https://stacks.math.columbia.edu/tag/0AME}{lemma contained in radical (Tag 0AME)}.
\end{proof}


\begin{lemma}
\label{lemma-existence-projective}
Let $A$ be Noetherian ring complete with respect to an ideal $I$.
Let $f : X \to \Spec(A)$ be a projective morphism. Let
$\mathcal{I} = I\mathcal{O}_X$.
Then the functor (\ref{equation-completion-functor}) is an equivalence.
\end{lemma}

\begin{proof}
We have already seen that (\ref{equation-completion-functor}) is
fully faithful in Lemma \ref{lemma-fully-faithful}. Thus it suffices
to show that the functor is essentially surjective.

\medskip\noindent
We first show that every object $(\mathcal{F}_n)$ of
$\textit{Coh}(X, \mathcal{I})$ is the quotient of an object
in the image of (\ref{equation-completion-functor}). 
Let $\mathcal{L}$ be an $f$-ample invertible sheaf on $X$.
Choose $d_0$ as in Lemma \ref{lemma-vanishing-projective}.
Choose a $d \geq d_0$ such that
$\mathcal{F}_1 \otimes \mathcal{L}^{\otimes d}$
is globally generated by some sections $s_{1, 1}, \ldots, s_{t, 1}$.
Since the transition maps of the system
$$
H^0(X, \mathcal{F}_{n + 1} \otimes \mathcal{L}^{\otimes d})
\longrightarrow
H^0(X, \mathcal{F}_n \otimes \mathcal{L}^{\otimes d})
$$
are surjective by the vanishing of $H^1$ we can lift
$s_{1, 1}, \ldots, s_{t, 1}$ to a compatible system of global sections
$s_{1, n}, \ldots, s_{t, n}$ of
$\mathcal{F}_n \otimes \mathcal{L}^{\otimes d}$.
These determine a compatible system of maps
$$
(s_{1, n}, \ldots, s_{t, n}) :
(\mathcal{L}^{\otimes -d})^{\oplus t} \longrightarrow \mathcal{F}_n
$$
Using Lemma \ref{lemma-inverse-systems-surjective}
we deduce that we have a surjective map
$$
\left((\mathcal{L}^{\otimes -d})^{\oplus t}\right)^\wedge
\longrightarrow
(\mathcal{F}_n)
$$
as desired.

\medskip\noindent
The result of the previous paragraph and the fact that
$\textit{Coh}(X, \mathcal{I})$ is abelian
(Lemma \ref{lemma-inverse-systems-abelian})
implies that
every object of $\textit{Coh}(X, \mathcal{I})$ is a cokernel
of a map between objects coming from $\textit{Coh}(\mathcal{O}_X)$.
As (\ref{equation-completion-functor}) is fully faithful and exact by
Lemmas \ref{lemma-fully-faithful} and \ref{lemma-exact}
we conclude.
\end{proof}
\end{document}
